% =====================================================================
%  IdeaTheory monograph --- shared preamble
%  PDF target only; no exotic packages, EPUB-friendly subset.
% =====================================================================
\usepackage[utf8]{inputenc}
\usepackage[T1]{fontenc}
\usepackage{lmodern}
\usepackage{microtype}
\usepackage[english]{babel}
\usepackage[a4paper,margin=1in]{geometry}
\usepackage{setspace}
\usepackage{parskip}

\usepackage{amsmath,amssymb,amsthm,mathtools}
\usepackage{stmaryrd}
\usepackage{mathrsfs}
\usepackage{bm}

\usepackage{graphicx}
\usepackage{xcolor}
\usepackage{array}
\usepackage{booktabs}
\usepackage{longtable}
\usepackage{tabularx}
\usepackage{multirow}
\usepackage{enumitem}
\usepackage{epigraph}
\usepackage{csquotes}

\usepackage[colorlinks=true,
  linkcolor=blue!55!black,
  citecolor=teal!60!black,
  urlcolor=blue!55!black]{hyperref}
\usepackage{cleveref}

%% ---------- Theorem environments -------------------------------------
\theoremstyle{plain}
\newtheorem{theorem}{Theorem}[chapter]
\newtheorem{lemma}[theorem]{Lemma}
\newtheorem{proposition}[theorem]{Proposition}
\newtheorem{corollary}[theorem]{Corollary}

\theoremstyle{definition}
\newtheorem{definition}[theorem]{Definition}
\newtheorem{example}[theorem]{Example}
\newtheorem{construction}[theorem]{Construction}
\newtheorem{notation}[theorem]{Notation}
\newtheorem{axiom}[theorem]{Axiom}

\theoremstyle{remark}
\newtheorem{remark}[theorem]{Remark}
\newtheorem{interpretation}[theorem]{Interpretation}
\newtheorem*{zooentry*}{Zoo Entry}

%% ---------- IdeaTheory notation --------------------------------------
\newcommand{\Ideas}{\mathcal{I}}
\newcommand{\IdeaAlg}{\mathbb{I}}
\newcommand{\compose}{\mathbin{\diamond}}
\newcommand{\unit}{\mathbf{1}}
\newcommand{\IdMon}{(\Ideas,\compose,\unit)}

\newcommand{\res}[2]{\langle #1 \mathbin{\Vert} #2 \rangle}
\newcommand{\Res}{\mathsf{R}}

\newcommand{\Auf}{\mathfrak{A}}
\newcommand{\preserve}{\sigma}
\newcommand{\negate}{\nu}
\newcommand{\elevate}{\eta}

\newcommand{\emerg}{\omega}
\newcommand{\dcoc}{\partial}
\newcommand{\curv}{\kappa}

\newcommand{\conj}[2]{{}^{#1}\!#2}
\newcommand{\chain}[1]{\mathsf{Ch}(#1)}
\newcommand{\equivIdea}{\sim_{\Ideas}}

\newcommand{\NN}{\mathbb{N}}
\newcommand{\ZZ}{\mathbb{Z}}
\newcommand{\QQ}{\mathbb{Q}}
\newcommand{\RR}{\mathbb{R}}
\newcommand{\CC}{\mathbb{C}}
\newcommand{\eqdef}{\mathrel{\vcentcolon=}}
\newcommand{\defeq}{\mathrel{\vcentcolon=}}

%% ---------- Zoo macros -----------------------------------------------
% Each "zoo entry" formalizes one informal theory of ideas.
% Usage:
%   \begin{zooentry}{Hegel: Aufhebung}{Hegel1812}
%      Original claim ... \par
%      \textbf{Formalization.} ... \par
%      \textbf{Status.} Theorem~3.4 ...
%   \end{zooentry}
\newenvironment{zooentry}[2]{%
  \par\medskip\noindent\textbf{Zoo entry: #1}\\
  \emph{Source:} \cite{#2}.\par\smallskip
}{\par\medskip}

\newcommand{\thinker}[1]{\textbf{#1}}
\newcommand{\formalclaim}{\par\textbf{Formal restatement.}\ }
\newcommand{\formalstatus}{\par\textbf{Status in Idea Theory.}\ }
\newcommand{\litcomment}{\par\textbf{Commentary.}\ }

% Chapter 2 (coalitions) macros
\newcommand{\aggregate}{\mathsf{aggregate}}
\newcommand{\Coalition}{\mathsf{Coalition}}
\newcommand{\insertAt}{\mathsf{insertAt}}
\newcommand{\chair}{\mathsf{chair}}
\newcommand{\members}{\mathsf{members}}
\newcommand{\collective}{\mathsf{collective}}
\newcommand{\Comm}{\mathfrak{C}}
\newcommand{\Asm}{\mathfrak{A}_{\mathrm{sm}}}
\newcommand{\List}{\mathsf{List}}
\newcommand{\replicate}{\mathsf{replicate}}
\newcommand{\xs}{\mathbf{x}}
\newcommand{\ys}{\mathbf{y}}
\newcommand{\zs}{\mathbf{z}}
